\chapter{Field Extensions and Galois Theory}\label{ch:b3:galois}

Can every equation be solved by radicals, as the quadratic formula
and Cardano's cubic formulas suggest? Can one trisect an angle
with ruler and compass? Both questions, open for centuries, are
answered --- negatively --- by a single idea of \'Evariste Galois:
attach to every polynomial a finite \emph{group} of symmetries of
its roots, and read the answer off the group. This chapter builds
the dictionary: field extensions and degrees, splitting fields and
algebraic closures, finite fields (a complete theory --- and the
promised cyclicity of $\mathbb F_q^\times$), separability, then
the Galois correspondence itself, with full proofs. We harvest:
the impossibility of the classical constructions, the structure of
cyclotomic fields, and the unsolvability of the quintic by
radicals --- \cref{ch:b3:groups}'s simplicity of $A_5$ striking
its target.

\section{Extensions, degree, algebraicity}

\begin{definition}\label{def:b3:galois:extension}
A \emph{field extension}\index{field extension} $L/K$ is a field
$L$ containing $K$ as a subfield; $L$ is then a $K$-vector space,
and the \emph{degree}\index{degree of an extension} $[L:K]$ is
its dimension. The extension is \emph{finite} if $[L:K] <
\infty$. The \emph{characteristic} of a field is the generator
$\geq 0$ of the kernel of $\Z \to K$, $n \mapsto n\cdot 1$: it is
$0$ or a prime $p$; correspondingly $K$ contains a smallest
subfield (\emph{prime field}) isomorphic to $\Q$ or to $\mathbb
F_p = \Z/p\Z$.
\end{definition}

\begin{theorem}[Tower law]\label{thm:b3:galois:tower}
If $K \subseteq L \subseteq M$, then $[M:K] = [M:L]\,[L:K]$: if
$(e_i)$ is a basis of $L$ over $K$ and $(f_j)$ a basis of $M$
over $L$, then $(e_if_j)$ is a basis of $M$ over
$K$.\index{tower law}
\end{theorem}

\begin{proof}
Generating: $x \in M$ writes $x = \sum_j \lambda_jf_j$
($\lambda_j \in L$), each $\lambda_j = \sum_i \mu_{ij}e_i$
($\mu_{ij} \in K$): $x = \sum_{i,j}\mu_{ij}e_if_j$. Independent:
$\sum_{i,j}\mu_{ij}e_if_j = 0$ rewrites $\sum_j (\sum_i
\mu_{ij}e_i)f_j = 0$; the inner sums are in $L$, so vanish
($f_j$ independent over $L$); then all $\mu_{ij} = 0$ ($e_i$
independent over $K$).
\end{proof}

\begin{definition}\label{def:b3:galois:algebraic}
Let $L/K$ and $\alpha \in L$. If some nonzero $P \in K[X]$ has
$P(\alpha) = 0$, $\alpha$ is \emph{algebraic}\index{algebraic
element} over $K$; the monic generator $\pi_\alpha$ of the ideal
$\{P : P(\alpha) = 0\}$ of $K[X]$ is its \emph{minimal
polynomial}\index{minimal polynomial (field theory)}, an
irreducible polynomial ($\pi = QR$ with $Q(\alpha) = 0$ forces $R$
constant by minimality of the degree). Otherwise $\alpha$ is
\emph{transcendental}. We write $K(\alpha)$ for the smallest
subfield of $L$ containing $K$ and $\alpha$, and $K[\alpha]$ for
the smallest subring.
\end{definition}

\begin{theorem}\label{thm:b3:galois:simple}
If $\alpha$ is algebraic over $K$ with $d = \deg\pi_\alpha$,
then
\[
K(\alpha) = K[\alpha] \cong K[X]/(\pi_\alpha),
\qquad [K(\alpha) : K] = d,
\]
with basis $1, \alpha, \dots, \alpha^{d-1}$. Conversely, if
$[L:K] < \infty$, every $\alpha \in L$ is algebraic of degree
$\deg \pi_\alpha$ dividing $[L:K]$.
\end{theorem}

\begin{proof}
Evaluation $K[X] \to L$, $P \mapsto P(\alpha)$, has image
$K[\alpha]$ and kernel $(\pi_\alpha)$; since $\pi_\alpha$ is
irreducible, $K[X]/(\pi_\alpha)$ is a \emph{field}
(\cref{prop:b3:rings:primemaximal}: in the PID $K[X]$,
irreducible generates a maximal ideal), so $K[\alpha]$ is a field
containing $K$ and $\alpha$: it equals $K(\alpha)$. The classes
of $1, X, \dots, X^{d-1}$ form a basis of the quotient (Euclidean
division), whence the basis and the degree. Conversely if $[L:K]
= n < \infty$: $1, \alpha, \dots, \alpha^n$ are dependent, giving
an annihilating polynomial; then $[K(\alpha):K] =
\deg\pi_\alpha$ divides $n$ by the tower law.
\end{proof}

\begin{corollary}\label{cor:b3:galois:algclosed}
If $\alpha, \beta$ are algebraic over $K$, so are $\alpha \pm
\beta$, $\alpha\beta$, $\alpha/\beta$ ($\beta \ne 0$): the
elements of $L$ algebraic over $K$ form a subfield of $L$.
Moreover algebraicity is transitive: algebraic over algebraic is
algebraic.
\end{corollary}

\begin{proof}
$K(\alpha, \beta) = (K(\alpha))(\beta)$ is finite over
$K(\alpha)$ ($\beta$ algebraic over $K \subseteq K(\alpha)$) and
$K(\alpha)/K$ is finite: by the tower law $[K(\alpha,\beta):K] <
\infty$, and every element of $K(\alpha, \beta)$ --- including
the four listed --- is algebraic (\cref{thm:b3:galois:simple}).
Transitivity: if $\beta$ is algebraic over $L$ and $L/K$ is
algebraic, the coefficients $c_0, \dots, c_{m-1}$ of
$\pi_{\beta/L}$ generate a finite extension $F = K(c_0, \dots,
c_{m-1})$ of $K$ (repeated tower law), and $F(\beta)/F$ is
finite: $[F(\beta):K] < \infty$, so $\beta$ is algebraic over
$K$.
\end{proof}

\begin{example}\label{ex:b3:galois:degrees}
$[\Q(\sqrt2):\Q] = 2$, $[\Q(\sqrt[3]2):\Q] = 3$ ($X^3 - 2$ is
irreducible: Eisenstein), $[\Q(\zeta_p):\Q] = p - 1$ for
$\zeta_p = \eu^{2\iu\pi/p}$ ($\Phi_p$ is irreducible,
\cref{ex:b3:rings:eisenstein}). The tower law is already a
weapon: $\sqrt[3]2 \notin \Q(\sqrt2)$, since $3 \nmid 2$.
\end{example}

\section{Splitting fields; algebraic closure}

\begin{theorem}[Splitting fields]\label{thm:b3:galois:splitting}
Let $P \in K[X]$ be nonconstant. There exists a \emph{splitting
field}\index{splitting field} of $P$ over $K$: an extension $L =
K(\alpha_1, \dots, \alpha_n)$ generated by roots of $P$ in which
$P$ splits into linear factors. It is unique up to
$K$-isomorphism, and $[L:K] \leq (\deg P)!$.
\end{theorem}

\begin{proof}
\emph{Existence}, by induction on $\deg P$: pick an irreducible
factor $Q$ of $P$; the field $K_1 = K[X]/(Q)$ contains the root
$\alpha_1 = \bar X$ of $Q$, hence of $P$; write $P = (X -
\alpha_1)P_1$ over $K_1$ and apply induction to $P_1$ over $K_1$;
degrees multiply to at most $n(n-1)\cdots = n!$.

\emph{Uniqueness} follows from the stronger \emph{isomorphism
extension lemma}: let $\sigma \colon K \to K'$ be an isomorphism,
$P \in K[X]$, $P^\sigma$ the polynomial with mapped coefficients,
$L, L'$ splitting fields of $P, P^\sigma$; then $\sigma$ extends
to an isomorphism $L \to L'$. Induction on $[L:K]$: if $P$ splits
in $K$, then $L = K$, and $L' = K'$ ($P^\sigma$ splits in $K'$,
and $L'$ is generated by its roots). Otherwise choose a root
$\alpha \in L \setminus K$ of an irreducible factor $Q$ of $P$
with $\deg Q \geq 2$; $Q^\sigma$ is an irreducible factor of
$P^\sigma$, with a root $\beta \in L'$; then
\[
K(\alpha) \cong K[X]/(Q) \xrightarrow{\ \sigma\ }
K'[X]/(Q^\sigma) \cong K'(\beta)
\]
extends $\sigma$ with $\alpha \mapsto \beta$. Now $L$ is a
splitting field of $P$ over $K(\alpha)$, and $L'$ of $P^\sigma$
over $K'(\beta)$, with $[L : K(\alpha)] < [L:K]$: induction
extends further to $L \to L'$.
\end{proof}

\begin{definition}\label{def:b3:galois:closure}
A field $\Omega$ is \emph{algebraically closed}\index{algebraically
closed field} if every nonconstant polynomial of $\Omega[X]$ has
a root in $\Omega$ (hence splits). An \emph{algebraic
closure}\index{algebraic closure} of $K$ is an algebraic
extension $\bar K/K$ with $\bar K$ algebraically closed.
\end{definition}

\begin{theorem}[Steinitz]\label{thm:b3:galois:steinitz}
Every field $K$ has an algebraic closure, unique up to
$K$-isomorphism.
\end{theorem}

\begin{proof}
\emph{Existence (Artin's construction).} Let $R = K[(X_f)_f]$ be
the polynomial ring with one variable $X_f$ per nonconstant
monic $f \in K[X]$, and $I$ the ideal generated by all
$f(X_f)$. $I$ is proper: a relation $1 = \sum_{i=1}^r g_i\,
f_i(X_{f_i})$ involves finitely many polynomials; in a common
splitting field $E$ of $f_1\cdots f_r$ pick roots $\alpha_i$ of
$f_i$ and evaluate $X_{f_i} \mapsto \alpha_i$ (other variables
$\mapsto 0$): $1 = 0$, absurd. Let $\mathfrak m \supseteq I$ be
maximal (\cref{thm:b3:rings:krull}; Zorn) and $K_1 = R/\mathfrak
m$: a field extension of $K$ in which every nonconstant $f \in
K[X]$ has a root, namely $\bar X_f$, and which is algebraic over
$K$ (it is generated by the $\bar X_f$, each algebraic).
Iterate: $K \subseteq K_1 \subseteq K_2 \subseteq \cdots$, where
$K_{n+1}$ does to $K_n$ what $K_1$ did to $K$, and let $\Omega =
\bigcup_n K_n$, a field. Any nonconstant $g \in \Omega[X]$ has
its finitely many coefficients in some $K_n$; an irreducible
factor of $g$ over $K_n$ has a root in $K_{n+1} \subseteq
\Omega$: $\Omega$ is algebraically closed, and algebraic over
$K$ (each $K_n$ is, by transitivity,
\cref{cor:b3:galois:algclosed}): $\Omega$ is an algebraic
closure.

\emph{Uniqueness.} Let $\Omega, \Omega'$ be two algebraic
closures. Consider the set of pairs $(E, \tau)$ where $K
\subseteq E \subseteq \Omega$ and $\tau\colon E \to \Omega'$ is a
$K$-embedding, ordered by extension; it is nonempty ($(K,
\mathrm{id})$) and inductive (union of a chain), so Zorn gives a
maximal $(E_0, \tau_0)$. If $E_0 \neq \Omega$, pick $\alpha \in
\Omega\setminus E_0$: $\pi_{\alpha/E_0}$ maps to a polynomial
over $\tau_0(E_0)$ having a root $\beta$ in the algebraically
closed $\Omega'$, and $\tau_0$ extends to $E_0(\alpha) \to
\Omega'$ ($\alpha \mapsto \beta$), contradicting maximality. So
there is a $K$-embedding $\tau \colon \Omega \to \Omega'$; its
image, isomorphic to $\Omega$, is algebraically closed, and
$\Omega'$ is algebraic over it: for $x \in \Omega'$,
$\pi_{x/\tau(\Omega)}$ splits over $\tau(\Omega)$, so $x \in
\tau(\Omega)$. Thus $\tau$ is onto: an isomorphism.
\end{proof}

\begin{remark}\label{rem:b3:galois:closureexamples}
For $K = \Q$ one may avoid the transfinite machinery: the
algebraic numbers $\bar\Q = \{z \in \C : z \text{ algebraic over
} \Q\}$ form an algebraic closure --- a subfield of $\C$ by
\cref{cor:b3:galois:algclosed}, algebraically closed because $\C$
is (d'Alembert--Gauss, proved by complex analysis in
\cref{ch:b3:holomorphic}) and roots of polynomials over $\bar\Q$
are algebraic over $\Q$ by transitivity.
\end{remark}

\section{Finite fields}

\begin{theorem}\label{thm:b3:galois:finitefields}
Let $p$ be prime, $n \geq 1$, $q = p^n$.
\begin{enumerate}
  \item A finite field has cardinality a prime power, and for
        each $q$ there is exactly one field $\mathbb F_q$ with
        $q$ elements up to isomorphism: the splitting field of
        $X^q - X$ over $\mathbb F_p$.\index{finite field}
  \item The \emph{Frobenius}\index{Frobenius endomorphism} $F
        \colon x \mapsto x^p$ is an automorphism of $\mathbb
        F_q$, and the automorphism group of $\mathbb F_q$ is
        cyclic of order $n$, generated by $F$.
  \item $\mathbb F_{p^m}$ embeds in $\mathbb F_{p^n}$ iff $m \mid
        n$.
\end{enumerate}
\end{theorem}

\begin{proof}
(1) A finite field $E$ has characteristic $p > 0$ and is a
finite-dimensional $\mathbb F_p$-vector space: $\abs E = p^n$.
Its multiplicative group has order $q - 1$, so every $x \in E$
satisfies $x^q = x$: $E$ consists of $q$ roots of $X^q - X$,
hence is a splitting field of it over $\mathbb F_p$ ---
determining $E$ up to isomorphism
(\cref{thm:b3:galois:splitting}). Conversely, in a splitting
field $L$ of $X^q - X$, the set $E$ of its roots is a
\emph{subfield}: $(x + y)^q = x^q + y^q$ by iterating the
freshman's dream $(a+b)^p = a^p + b^p$ ($p \mid \binom pk$), and
$(xy)^q = x^qy^q$, $(x^{-1})^q = (x^q)^{-1}$; it has exactly $q$
elements since $X^q - X$ is separable: its derivative is
$qX^{q-1} - 1 = -1$ (as $p \mid q$), coprime to it, so no
repeated roots. Thus $L = E$ has $q$ elements.

(2) $F$ is a field morphism (freshman's dream), injective
(fields), hence bijective on the finite $\mathbb F_q$. $F^n =
\mathrm{id}$ ($x^q = x$), and no smaller power is the identity:
$F^m = \mathrm{id}$ means all $q$ elements are roots of
$X^{p^m} - X$, forcing $p^m \geq q$. So $\langle F\rangle$ is
cyclic of order $n$; and there are no other automorphisms, by
the bound $\abs{\operatorname{Aut}} \leq [\,\mathbb F_q :
\mathbb F_p\,] = n$ proved below
(\cref{prop:b3:galois:embeddings} with $L = \mathbb F_q$, $K =
\mathbb F_p$: automorphisms fix the prime field).

(3) If $\mathbb F_{p^m} \subseteq \mathbb F_{p^n}$, the tower law
gives $p^n = (p^m)^d$: $m \mid n$. Conversely if $m \mid n$, then
$p^m - 1 \mid p^n - 1$ (geometric sum), so $X^{p^m} - X$ divides
$X^{p^n} - X$ (same argument on exponents: $X^{a} - 1 \mid X^{b}
- 1$ when $a \mid b$), and the roots of the former inside
$\mathbb F_{p^n}$ form the required subfield, of cardinality
$p^m$ (separability as in (1)).
\end{proof}

\begin{theorem}[Cyclicity]\label{thm:b3:galois:cyclic}
Every finite subgroup of the multiplicative group of a field is
cyclic. In particular $\mathbb F_q^\times \cong
\Z/(q-1)\Z$.\index{multiplicative group of a finite field}
\end{theorem}

\begin{proof}
Let $G \leq K^\times$ be finite. By the structure theorem
(\cref{cor:b3:modules:abelian}), $G \cong \Z/d_1 \times \dots
\times \Z/d_s$ with $d_1 \mid \dots \mid d_s$. Every $x \in G$
then satisfies $x^{d_s} = 1$; but $X^{d_s} - 1$ has at most
$d_s$ roots in the field $K$: $\abs G = d_1\cdots d_s \leq d_s$,
forcing $s = 1$: $G$ is cyclic.
\end{proof}

\begin{example}\label{ex:b3:galois:f8}
$\mathbb F_8 = \mathbb F_2[X]/(X^3 + X + 1)$: the cubic has no
root in $\mathbb F_2$, hence is irreducible. Writing $\omega =
\bar X$: $\mathbb F_8^\times$ is cyclic of order $7$, so
\emph{every} element $\neq 0, 1$ generates. The subfields of
$\mathbb F_{p^{12}}$ form the divisor lattice of $12$:
$\mathbb F_p, \mathbb F_{p^2}, \mathbb F_{p^3}, \mathbb F_{p^4},
\mathbb F_{p^6}, \mathbb F_{p^{12}}$ --- a first, complete
instance of the Galois correspondence.
\end{example}

\section{Separability and embeddings}

\begin{definition}\label{def:b3:galois:separable}
A polynomial $P \in K[X]$ is \emph{separable}\index{separable
polynomial} if it has no repeated root in a splitting field ---
equivalently $\gcd(P, P') = 1$ (a repeated root is a common
root; conversely, over the splitting field, a common root is
repeated; and the gcd does not change under field extension,
\cref{cor:b3:modules:descent}'s argument). An algebraic element
is separable if its minimal polynomial is; an extension $L/K$ is
separable if all its elements are.
\end{definition}

\begin{proposition}\label{prop:b3:galois:perfect}
An \emph{irreducible} $P \in K[X]$ is separable unless $P' = 0$,
which forces $\operatorname{char} K = p > 0$ and $P \in K[X^p]$.
Consequently every algebraic extension of a field of
characteristic $0$, and of a finite field, is separable (such
fields are called \emph{perfect}\index{perfect field}).
\end{proposition}

\begin{proof}
$\gcd(P, P')$ divides $P$; if it is not $1$, irreducibility
forces $\gcd = P$ (up to a constant), so $P \mid P'$ with $\deg
P' < \deg P$: $P' = 0$. Writing $P = \sum a_kX^k$: $ka_k = 0$
for all $k$, so in characteristic $0$, $P$ is constant (excluded);
in characteristic $p$, $a_k = 0$ unless $p \mid k$: $P =
Q(X^p)$. Over a finite field, every element is a $p$-th power
(Frobenius is onto), so $Q(X^p) = \sum b_k^p X^{pk} = (\sum
b_kX^k)^p$ is not irreducible: $P' = 0$ cannot happen for
irreducible $P$ there either.
\end{proof}

\begin{proposition}[Counting embeddings]\label{prop:b3:galois:embeddings}
Let $L = K(\alpha_1, \dots, \alpha_r)$ be finite over $K$, and
$\sigma \colon K \to \Omega$ an embedding into an algebraically
closed field. Then the number of extensions of $\sigma$ to $L$
is at most $[L:K]$, with equality if $L/K$ is separable. In
particular $\abs{\operatorname{Aut}_K(L)} \leq [L:K]$.
\end{proposition}

\begin{proof}
Induction on $[L:K]$ via simple steps. For $L = K(\alpha)$: an
extension $\tau$ is determined by $\tau(\alpha)$, which must be a
root in $\Omega$ of $\pi_\alpha^\sigma$; conversely each such
root gives one extension ($K(\alpha) \cong K[X]/(\pi_\alpha)$).
The number of extensions is the number of \emph{distinct} roots
of $\pi_\alpha^\sigma$ in $\Omega$: at most $\deg\pi_\alpha =
[K(\alpha):K]$, with equality iff $\pi_\alpha$ is separable
(separability of $\pi^\sigma$ and $\pi$ agree: gcd with the
derivative is preserved by $\sigma$). In general, factor $L =
K(\alpha_1)(\alpha_2, \dots)$: extensions of $\sigma$ to
$K(\alpha_1)$ number $\leq [K(\alpha_1):K]$, and each extends in
$\leq [L : K(\alpha_1)]$ ways by induction; multiply (tower
law). In the separable case both counts are equalities: minimal
polynomials over the bigger field $K(\alpha_1)$ divide those
over $K$, hence remain separable.
\end{proof}

\begin{theorem}[Primitive element]\label{thm:b3:galois:primitive}
Every finite \emph{separable} extension is simple: $L =
K(\gamma)$ for some $\gamma$.\index{primitive element theorem}
\end{theorem}

\begin{proof}
If $K$ is finite, so is $L$, and a generator $\gamma$ of the
cyclic group $L^\times$ (\cref{thm:b3:galois:cyclic}) does it.
Let $K$ be infinite; by induction it suffices to treat $L =
K(\alpha, \beta)$. Let $n = [L:K]$; by
\cref{prop:b3:galois:embeddings} there are $n$ distinct
$K$-embeddings $\sigma_1, \dots, \sigma_n \colon L \to \Omega$
($\Omega$ an algebraic closure). The polynomial
\[
D(T)=\prod_{i<j}\bigl[\bigl(\sigma_i(\alpha)-\sigma_j(\alpha)
\bigr) + T\bigl(\sigma_i(\beta) - \sigma_j(\beta)\bigr)\bigr]
\]
is not identically zero: a factor vanishes identically only if
$\sigma_i, \sigma_j$ agree on both $\alpha$ and $\beta$, hence on
$L$ --- excluded for $i \neq j$. As $K$ is infinite, pick $c \in
K$ with $D(c) \ne 0$: then the $n$ elements $\sigma_i(\alpha +
c\beta)$ are pairwise distinct, so $\gamma = \alpha + c\beta$
has at least $n$ distinct conjugates in $\Omega$, i.e.\
$\deg \pi_\gamma \geq n$: $[K(\gamma):K] \geq n = [L:K]$ forces
$L = K(\gamma)$.
\end{proof}

\section{The Galois correspondence}

\begin{definition}\label{def:b3:galois:galois}
A finite extension $L/K$ is \emph{Galois}\index{Galois extension}
if it is the splitting field of a \emph{separable} polynomial
over $K$. Its \emph{Galois group}\index{Galois group} is
$\operatorname{Gal}(L/K) = \operatorname{Aut}_K(L)$, the group of
field automorphisms of $L$ fixing $K$ pointwise.
\end{definition}

\begin{proposition}\label{prop:b3:galois:galoisorder}
If $L/K$ is Galois, then $\abs{\operatorname{Gal}(L/K)} =
[L:K]$; moreover $L/F$ is Galois for every intermediate field $K
\subseteq F \subseteq L$, and every $F$-embedding $L \to \Omega
\supseteq L$ has image $L$ (\emph{normality}).
\end{proposition}

\begin{proof}
Let $L$ split the separable $P$ over $K$, and fix an algebraic
closure $\Omega \supseteq L$. $L/K$ is separable: it is generated
by roots of $P$; separability of every element follows from the
equality case below, but let us argue directly ---
\cref{prop:b3:galois:embeddings} applied to the generators (roots
of the separable $P$, whose minimal polynomials divide $P$)
yields exactly $[L:K]$ extensions of $K \hookrightarrow \Omega$
(in the inductive step, the minimal polynomial of a root of $P$
over an intermediate field still divides $P$, hence is
separable). Each such embedding $\tau \colon L \to \Omega$
permutes the roots of $P$ ($\tau$ fixes the coefficients), and
$L$ is generated by them: $\tau(L) = L$. Hence embeddings $=$
automorphisms: $\abs{\operatorname{Gal}(L/K)} = [L:K]$. For
intermediate $F$: $L$ is also the splitting field of $P$ over
$F$, and $P$ remains separable: $L/F$ is Galois; the same
argument gives normality over $F$.
\end{proof}

\begin{lemma}[Artin]\label{lem:b3:galois:artin}
Let $G$ be a finite group of automorphisms of a field $L$ and
$K = L^G = \{x : \sigma(x) = x\ \forall\sigma \in G\}$ its fixed
field. Then $[L : L^G] \leq \abs G$.
\end{lemma}

\begin{proof}
Let $n = \abs G$, $G = \{\sigma_1, \dots, \sigma_n\}$, and
suppose $x_1, \dots, x_{n+1} \in L$ are linearly independent over
$K$. The homogeneous linear system of $n$ equations in $n+1$
unknowns $(c_j)$ over $L$,
\[
\sum_{j=1}^{n+1} c_j\,\sigma_i(x_j) = 0
\qquad (i = 1, \dots, n),
\]
has a nonzero solution; choose one with the \emph{fewest} nonzero
entries, say $c_1, \dots, c_r \neq 0$ (renumbering), $r \geq 2$
(a single $c_j\sigma_i(x_j) = 0$ is impossible), normalized
$c_r = 1$. Not all $c_j$ lie in $K$: the equation for $\sigma_i =
\mathrm{id}$ would contradict independence; say $c_1 \notin K$,
so $\tau(c_1) \ne c_1$ for some $\tau \in G$. Apply $\tau$ to all
equations: since $\tau\sigma_i$ runs over $G$, the vector
$(\tau(c_j))_j$ is another solution; subtracting, $(c_j -
\tau(c_j))_j$ is a solution with fewer nonzero entries (the
$r$-th entry $1 - 1 = 0$ vanishes, the first does not) and not
zero: contradiction. So any $n+1$ elements are dependent: $[L:K]
\leq n$.
\end{proof}

\begin{theorem}[Fundamental theorem of Galois theory]
\label{thm:b3:galois:fundamental}
Let $L/K$ be a Galois extension with group $G =
\operatorname{Gal}(L/K)$.\index{Galois correspondence}
\begin{enumerate}
  \item $L^G = K$.
  \item The maps $H \mapsto L^H$ and $F \mapsto
        \operatorname{Gal}(L/F)$ are mutually inverse,
        inclusion-reversing bijections between subgroups of $G$
        and intermediate fields $K \subseteq F \subseteq L$;
        moreover $[L : L^H] = \abs H$ and $[L^H : K] = [G : H]$.
  \item $H \trianglelefteq G$ iff $L^H/K$ is Galois, and then
        restriction induces $\operatorname{Gal}(L^H/K) \cong
        G/H$.
\end{enumerate}
\end{theorem}

\begin{proof}
(1) Clearly $K \subseteq L^G$. Conversely let $\alpha \in L
\setminus K$; we exhibit $\sigma \in G$ with $\sigma(\alpha) \ne
\alpha$. The minimal polynomial $\pi_\alpha$ over $K$ has degree
$\geq 2$ and is separable ($L/K$ separable,
\cref{prop:b3:galois:galoisorder}), so it has another root
$\beta \neq \alpha$ in an algebraic closure $\Omega \supseteq L$.
Extend the $K$-embedding $K(\alpha) \to \Omega$, $\alpha \mapsto
\beta$, to an embedding $\tau\colon L \to \Omega$
(\cref{prop:b3:galois:embeddings}); by normality
(\cref{prop:b3:galois:galoisorder}) $\tau(L) = L$, so $\tau \in
G$, $\beta = \tau(\alpha) \in L$, and $\tau(\alpha) \neq
\alpha$.

(2) For a subgroup $H$: $L/L^H$ is Galois
(\cref{prop:b3:galois:galoisorder}), and
$\operatorname{Gal}(L/L^H) \supseteq H$ trivially, so $[L:L^H] =
\abs{\operatorname{Gal}(L/L^H)} \geq \abs H$; Artin's lemma
gives $[L:L^H] \leq \abs H$: equality, and
$\operatorname{Gal}(L/L^H) = H$. For an intermediate field $F$:
$L/F$ Galois gives $L^{\operatorname{Gal}(L/F)} = F$ by (1)
applied to $L/F$. The two maps are mutually inverse; they
reverse inclusions evidently. Degrees: $[L:L^H] = \abs H$ just
proved, and $[L^H:K] = [L:K]/[L:L^H] = \abs G/\abs H$.

(3) For $\sigma \in G$ and $H \leq G$: $\sigma(L^H) =
L^{\sigma H\sigma^{-1}}$ (direct check). By the bijection,
$\sigma(L^H) = L^H$ for all $\sigma$ iff $H \trianglelefteq G$.
Now if $H \trianglelefteq G$, set $F = L^H$: every $\sigma \in
G$ restricts to an automorphism of $F$, giving a morphism $\rho
\colon G \to \operatorname{Aut}_K(F)$ with kernel $\{\sigma :
\sigma\restriction_F = \mathrm{id}\} = \operatorname{Gal}(L/F) =
H$. So $G/H$ embeds in $\operatorname{Aut}_K(F)$, whence
$\abs{\operatorname{Aut}_K(F)} \geq [G:H] = [F:K]$; the reverse
inequality always holds (\cref{prop:b3:galois:embeddings}):
$\abs{\operatorname{Aut}_K(F)} = [F:K]$ and $\rho$ is onto. It
remains to see $F/K$ is Galois: $F$ is separable over $K$
(inside the separable $L/K$), and $F = K(\gamma)$
(\cref{thm:b3:galois:primitive}); the polynomial
$\prod_{\sigma \in G/H}\bigl(X - \sigma(\gamma)\bigr)$ (product
over the distinct images, which lie in $F$: $\sigma(F) =
L^{\sigma H\sigma^{-1}} = L^H = F$ by normality of $H$) has
coefficients
fixed by $G$, hence in $K$ by (1): it is a separable polynomial
of $K[X]$ split by $F$, and its roots generate $F$: $F/K$ is
Galois. Conversely, if $F = L^H$ with $F/K$ Galois, normality of
$F$ (\cref{prop:b3:galois:galoisorder}, applied to embeddings
$F \to \Omega$ restricted from elements of $G$) gives
$\sigma(F) = F$ for all $\sigma \in G$, i.e.\ $H
\trianglelefteq G$.
\end{proof}

\begin{omfigure}
\begin{tikzpicture}[xscale=1.5, yscale=1.15]
  \tikzset{fd/.style={font=\small}}
  % subgroup lattice of S3 (left, upside down to mirror fields)
  \node[fd] (G)  at (-3.2, 0)   {$S_3$};
  \node[fd] (A)  at (-4.4, 1.2) {$\langle(1\,2\,3)\rangle$};
  \node[fd] (T1) at (-3.6, 1.2) {$\langle(1\,2)\rangle$};
  \node[fd] (T2) at (-2.8, 1.2) {$\langle(1\,3)\rangle$};
  \node[fd] (T3) at (-2.0, 1.2) {$\langle(2\,3)\rangle$};
  \node[fd] (E)  at (-3.2, 2.4) {$\{e\}$};
  \draw (G) -- (A); \draw (G) -- (T1); \draw (G) -- (T2);
  \draw (G) -- (T3);
  \draw (E) -- (A); \draw (E) -- (T1); \draw (E) -- (T2);
  \draw (E) -- (T3);
  % field lattice (right)
  \node[fd] (Q)  at (1.6, 0)   {$\Q$};
  \node[fd] (W)  at (0.4, 1.2) {$\Q(\iu\sqrt3)$};
  \node[fd] (R1) at (1.4, 1.2) {$\Q(\sqrt[3]2)$};
  \node[fd] (R2) at (2.4, 1.2) {$\Q(j\sqrt[3]2)$};
  \node[fd] (R3) at (3.5, 1.2) {$\Q(j^2\sqrt[3]2)$};
  \node[fd] (L)  at (1.6, 2.4) {$L = \Q(\sqrt[3]2, j)$};
  \draw (Q) -- (W); \draw (Q) -- (R1); \draw (Q) -- (R2);
  \draw (Q) -- (R3);
  \draw (L) -- (W); \draw (L) -- (R1); \draw (L) -- (R2);
  \draw (L) -- (R3);
  \draw[->, omDef, thick, shorten >=6pt, shorten <=6pt]
    (-1.4, 1.2) -- (-0.6, 1.2);
\end{tikzpicture}

{\small The Galois correspondence for the splitting field $L$ of
$X^3 - 2$ over $\Q$ ($j = \eu^{2\iu\pi/3}$): subgroups of
$\operatorname{Gal}(L/\Q) \cong S_3$ (left, order reversed)
match intermediate fields (right). The unique normal proper
subgroup $\langle(1\,2\,3)\rangle$ corresponds to the unique
subextension $\Q(\iu\sqrt3)/\Q$ that is Galois; the three
conjugate subgroups $\langle(i\,j)\rangle$ correspond to the
three conjugate cubic fields $\Q(j^k\sqrt[3]2)$, none of them
normal over $\Q$.}
\end{omfigure}

\section{Cyclotomic extensions}

\begin{definition}\label{def:b3:galois:cyclotomic}
Let $n \geq 1$ and $\zeta_n = \eu^{2\iu\pi/n}$. The $n$-th
\emph{cyclotomic polynomial}\index{cyclotomic polynomial} is
$\Phi_n = \prod_{\gcd(k,n)=1,\ 1 \le k \le n} \bigl(X -
\zeta_n^k\bigr)$, of degree $\varphi(n)$; grouping the roots of
$X^n - 1$ by exact order, $X^n - 1 = \prod_{d \mid n}\Phi_d$,
which shows inductively that $\Phi_n \in \Z[X]$ (Euclidean
division of monic integer polynomials).
\end{definition}

\begin{theorem}\label{thm:b3:galois:cyclotomicirred}
$\Phi_n$ is irreducible over $\Q$; hence $[\Q(\zeta_n) : \Q] =
\varphi(n)$ and
\[
\operatorname{Gal}\bigl(\Q(\zeta_n)/\Q\bigr) \;\cong\;
(\Z/n\Z)^\times,
\qquad \sigma_a(\zeta_n) = \zeta_n^a .
\]
The extension $\Q(\zeta_n)/\Q$ is thus Galois with \emph{abelian}
group.
\end{theorem}

\begin{proof}
Let $f = \pi_{\zeta_n}$, so $\Phi_n = fg$ with $f, g \in \Z[X]$
monic (Gauss's lemma \cref{lem:b3:rings:gauss}: contents
multiply, all polynomials monic). \emph{Claim: if $\zeta$ is a
root of $f$ and $p \nmid n$ is prime, then $\zeta^p$ is a root
of $f$.} Otherwise $\zeta^p$ is a root of $g$ (it is a primitive
$n$-th root of unity), so $\zeta$ is a root of $g(X^p)$, and $f
\mid g(X^p)$ in $\Z[X]$ (minimal polynomial, then Gauss again).
Reduce mod $p$: $\bar g(X^p) = \bar g(X)^p$ (Frobenius on
$\mathbb F_p[X]$: coefficientwise $a^p = a$, and freshman's
dream), so $\bar f \mid \bar g^{\,p}$: $\bar f$ and $\bar g$
share an irreducible factor, and $\bar\Phi_n = \bar f\bar g$ has
a repeated factor. Then so does $X^n - \bar 1$; but its
derivative $\bar nX^{n-1}$ is coprime to it ($p \nmid n$, and
$0$ is not a root): contradiction.

Every primitive root $\zeta_n^k$ ($\gcd(k, n) = 1$) is obtained
from $\zeta_n$ by successive prime powers not dividing $n$
(factor $k$): the claim propagates, so every primitive root is a
root of $f$: $f = \Phi_n$, irreducible. Consequently
$[\Q(\zeta_n):\Q] = \varphi(n)$, and $\Q(\zeta_n)$ is the
splitting field of the separable $X^n - 1$ (all roots are powers
of $\zeta_n$): Galois. An automorphism $\sigma$ sends $\zeta_n$
to another primitive root $\zeta_n^{a(\sigma)}$, and $\sigma
\mapsto a(\sigma)$ is an injective morphism into
$(\Z/n\Z)^\times$; both groups have order $\varphi(n)$:
isomorphism.
\end{proof}

\section{Ruler and compass}

\begin{definition}\label{def:b3:galois:constructible}
Identify the plane with $\C$; start from $\{0, 1\}$. A point is
\emph{constructible}\index{constructible number} if it is
obtainable by finitely many intersections of lines through two
already-constructed points and circles centered at a constructed
point with radius a distance of two constructed points.
\end{definition}

\begin{theorem}[Wantzel]\label{thm:b3:galois:wantzel}
$z \in \C$ is constructible iff there is a tower $\Q = F_0
\subseteq F_1 \subseteq \dots \subseteq F_r$ with $[F_{i+1} :
F_i] = 2$ and $z \in F_r$. In particular, a constructible number
is algebraic of degree a power of $2$ over
$\Q$.\index{Wantzel's theorem}
\end{theorem}

\begin{proof}
($\Rightarrow$) The coordinates of the intersection of two lines
through points with coordinates in a subfield $F \subseteq \R$
solve a linear system over $F$: they stay in $F$. Line--circle
and circle--circle intersections lead, after eliminating the
linear part (subtracting the two circle equations gives a line),
to a quadratic equation over $F$: the new coordinates lie in $F$
or in $F(\sqrt d)$ for some $d \in F$, $d > 0$. By induction,
every constructed point has coordinates in a tower of quadratic
extensions of $\Q$; and $z = x + \iu y$ is in a quadratic tower
too (adjoin $\iu$: one more quadratic step). The degree
consequence: $[\Q(z):\Q]$ divides $[F_r : \Q] = 2^r$ (tower
law).

($\Leftarrow$) The constructible numbers form a field: sums and
differences by parallelograms (parallels are constructible:
drop and raise perpendiculars twice --- the classical
perpendicular through a point uses one circle and two arcs);
products and quotients by Thales' intercept configurations
(given lengths $a, b$ construct $ab$ and $a/b$ with similar
triangles on two rays). And the field is closed under square
roots: for $a > 0$, the circle of diameter $1 + a$ and the
perpendicular at the junction point meet at height $\sqrt a$
(altitude-geometric-mean relation in a right triangle); for a
complex $w = \rho\eu^{\iu\theta}$, construct $\sqrt\rho$ and
bisect $\theta$ (angle bisection is a compass construction).
Real and imaginary parts of members of a quadratic tower are
therefore constructible by induction on the tower: each step
adjoins roots of a quadratic, expressible by field operations
and one square root of an already-constructed number (the
quadratic formula; in characteristic $0$).
\end{proof}

\begin{corollary}\label{cor:b3:galois:impossible}
The three classical problems are unsolvable by ruler and
compass:\index{ruler and compass}
\begin{enumerate}
  \item \emph{Duplication of the cube}: $\sqrt[3]2$ has degree
        $3$, not a power of $2$.
  \item \emph{Trisection of the angle}: trisecting $60^\circ$
        requires $\cos 20^\circ$, a root of the irreducible $8X^3
        - 6X - 1$: degree $3$.
  \item \emph{Squaring the circle}: $\sqrt\pi$ is transcendental
        ($\pi$ is --- Lindemann's theorem, admitted here: its
        proof belongs to a course in transcendence theory).
\end{enumerate}
Also, the regular $n$-gon is constructible iff
$\varphi(n)$ is a power of $2$ (Gauss--Wantzel; the ``if'' uses
the weekend problem's method, the ``only if'' is
\cref{thm:b3:galois:wantzel} applied to $\zeta_n$, of degree
$\varphi(n)$). For $n = 7$: $\varphi(7) = 6$: the regular
heptagon is impossible; for $n = 17$: $\varphi(17) = 16 = 2^4$:
constructible --- the weekend problem constructs it.
\end{corollary}

\begin{proof}
(1) $X^3 - 2$ is irreducible (Eisenstein). (2) From $\cos
3\theta = 4\cos^3\theta - 3\cos\theta$ with $3\theta =
60^\circ$: $8c^3 - 6c = 1$ for $c = \cos 20^\circ$; the cubic
$8X^3 - 6X - 1$ has no rational root (candidates $\pm1, \pm\frac
12, \pm\frac14, \pm\frac18$ fail), hence is irreducible: degree
$3$. A general $60^\circ$ angle is constructible, so a trisector
would construct $c$. (3) If $\sqrt\pi$ were constructible it
would be algebraic, hence also $\pi$. The $n$-gon statement: the
degree of $\zeta_n$ is $\varphi(n)$
(\cref{thm:b3:galois:cyclotomicirred}); necessity follows from
Wantzel; for sufficiency, the Galois group, abelian of order
$2^m$, admits a chain of index-$2$ subgroups (a finite $2$-group
does: \cref{exo:b3:groups:10}), whose fixed fields form a
quadratic tower ending at $\Q(\zeta_n)$
(\cref{thm:b3:galois:fundamental}); conclude by
\cref{thm:b3:galois:wantzel}.
\end{proof}

\section{Solvability by radicals}

\begin{definition}\label{def:b3:galois:radical}
An extension $L/K$ (characteristic $0$ throughout this section)
is \emph{radical} if there is a tower $K = F_0 \subseteq \dots
\subseteq F_r = L$ with $F_{i+1} = F_i(\alpha_i)$,
$\alpha_i^{n_i} \in F_i$: each step adjoins an $n_i$-th root. A
polynomial $P \in K[X]$ is \emph{solvable by
radicals}\index{solvable by radicals} if its splitting field is
contained in some radical extension of $K$.
\end{definition}

\begin{lemma}\label{lem:b3:galois:cyclicsteps}
Let $K$ contain a primitive $n$-th root of unity $\zeta$, i.e.\
$\zeta$ of order $n$ in $K^\times$, and $a \in K^\times$. Then
$K(\sqrt[n]a)/K$ is Galois with \emph{cyclic} group. Conversely
--- not needed below --- every cyclic extension of degree $n$ is
of this form. Moreover $K(\zeta_n)/K$ is Galois with
\emph{abelian} group, for any $K$ of characteristic $0$.
\end{lemma}

\begin{proof}
$X^n - a$ is separable ($\gcd$ with $nX^{n-1}$: $a \neq 0$) and
splits in $K(\alpha)$, $\alpha^n = a$: its roots are the
$\zeta^k\alpha \in K(\alpha)$. So $K(\alpha)/K$ is Galois; the
map $\sigma \mapsto \sigma(\alpha)/\alpha \in \mu_n = \langle
\zeta\rangle$ is an injective morphism ($\sigma\tau(\alpha) =
\sigma(\tau(\alpha)/\alpha \cdot \alpha) =
\tau(\alpha)/\alpha\cdot\sigma(\alpha)$, the quotient being in
$K$), into a cyclic group: $\operatorname{Gal}$ is cyclic. The
converse is Kummer theory, which we shall not need (see the
remark below). For $K(\zeta_n)$: it splits the separable $X^n -
1$, and $\sigma \mapsto a(\sigma)$ with $\sigma(\zeta_n) =
\zeta_n^{a(\sigma)}$ embeds the group in the abelian
$(\Z/n\Z)^\times$ as in \cref{thm:b3:galois:cyclotomicirred}
(injectivity only needs $\zeta_n$ to generate the roots of
unity involved).
\end{proof}

\begin{theorem}[Galois]\label{thm:b3:galois:solvable}
Let $K$ be of characteristic $0$ and $P \in K[X]$ with splitting
field $L$. If $P$ is solvable by radicals, then
$\operatorname{Gal}(L/K)$ is a solvable group. (The converse is
also true; we shall not need it.)\index{Galois' solvability
theorem}
\end{theorem}

\begin{proof}
Step 1: \emph{enlarge the radical tower to a Galois one.} Let $L
\subseteq M$ with $M/K$ radical, with root exponents $n_1, \dots,
n_r$ and $n = n_1\cdots n_r$. First adjoin $\zeta_n$: the tower
$K \subseteq K(\zeta_n) \subseteq M(\zeta_n)$ is still radical
($\zeta_n$ is a root of unity: a radical step, $\zeta_n^n = 1$),
and its steps beyond the first happen over fields containing the
needed roots of unity. Next, replace $M(\zeta_n)$ by the
composite $N$ of all $\sigma(M(\zeta_n))$, $\sigma$ ranging over
the (finitely many) $K$-embeddings of $M(\zeta_n)$ into a fixed
algebraic closure: $N$ is the splitting field of the product of
minimal polynomials of a generating set (characteristic $0$:
finite and separable), hence $N/K$ is Galois; and $N$ is
radical over $K$: each $\sigma(M(\zeta_n))$ is radical over $K$
(apply $\sigma$ to a radical tower), and a composite of radical
extensions is radical (concatenate the towers: if $F'/K$ is
radical with tower adjoining $\beta_j$, then $F''(\beta_j)$-type
steps remain radical over any bigger base).

Step 2: \emph{read solvability off the Galois tower.} So assume
$L \subseteq N$, $N/K$ Galois and radical with tower $K
\subseteq K(\zeta_n) = E_0 \subseteq E_1 \subseteq \dots
\subseteq E_s = N$, each $E_{i+1} = E_i(\sqrt[n_i]{a_i})$ with
$\zeta_{n_i} \in E_0 \subseteq E_i$. Let $G =
\operatorname{Gal}(N/K)$ and $G_i = \operatorname{Gal}(N/E_i)$:
a decreasing chain $G \supseteq G_0 \supseteq G_1 \supseteq
\dots \supseteq G_s = \{e\}$. Each $E_{i+1}/E_i$ is Galois with
cyclic group (\cref{lem:b3:galois:cyclicsteps}), so by the
fundamental theorem applied to the Galois extension $N/E_i$
(\cref{thm:b3:galois:fundamental}(3), with ambient group $G_i$):
$G_{i+1} \trianglelefteq G_i$ with $G_i/G_{i+1} \cong
\operatorname{Gal}(E_{i+1}/E_i)$ cyclic. Similarly $E_0/K$ is
Galois with abelian group $G/G_0$
(\cref{lem:b3:galois:cyclicsteps}). The chain exhibits $G$ as
solvable (\cref{prop:b3:groups:derived}). Finally
$\operatorname{Gal}(L/K)$ is a \emph{quotient} of $G$: $L/K$ is
Galois ($P$ separable in characteristic $0$) and restriction $G
\to \operatorname{Gal}(L/K)$ is onto
(\cref{thm:b3:galois:fundamental}(3) with $H =
\operatorname{Gal}(N/L)$); quotients of solvable groups are
solvable.
\end{proof}

\begin{corollary}[Unsolvability of the quintic]
\label{cor:b3:galois:quintic}
There are polynomials of degree $5$ over $\Q$ that are not
solvable by radicals: for instance $X^5 - 4X + 2$, whose Galois
group is $S_5$ (\cref{exo:b3:galois:11}), a non-solvable group
(\cref{cor:b3:groups:snnotsolvable}). No general formula in
radicals can exist for degree $\geq 5$.\index{quintic,
unsolvability}
\end{corollary}

\begin{remark}\label{rem:b3:galois:converse}
The converse of \cref{thm:b3:galois:solvable} --- a solvable
Galois group implies solvability by radicals --- is proved by
descending the derived series and showing every cyclic extension
(with enough roots of unity) is radical, via \emph{Lagrange
resolvents}; it explains why degrees $2, 3, 4$ have formulas:
$S_2, S_3, S_4$ are solvable
(\cref{ex:b3:groups:solvableexamples}). We leave it admitted at
this level; a full treatment belongs to a master's course, but
\cref{exo:b3:galois:8} makes it concrete for the cubic.
\end{remark}

\section{Exercises}

\begin{exercise}[$\star$]\label{exo:b3:galois:1}
Show $[\Q(\sqrt2, \sqrt3):\Q] = 4$, that $\Q(\sqrt2 + \sqrt3) =
\Q(\sqrt2, \sqrt3)$, and compute the minimal polynomial of
$\sqrt2 + \sqrt3$ over $\Q$.
\end{exercise}

\begin{exercise}[$\star$]\label{exo:b3:galois:2}
Let $\alpha = \sqrt[3]2$. Show that $\Q(\alpha)/\Q$ is not
normal (exhibit an embedding $\Q(\alpha) \to \C$ whose image is
not $\Q(\alpha)$), determine the splitting field $L$ of $X^3 -
2$ and $[L:\Q]$, and check $\operatorname{Aut}_\Q(\Q(\alpha)) =
\{\mathrm{id}\}$: for non-Galois extensions, the automorphism
group can be far smaller than the degree.
\end{exercise}

\begin{exercise}[$\star$]\label{exo:b3:galois:3}
Construct $\mathbb F_9$ as $\mathbb F_3[X]/(X^2+1)$ and find a
generator of $\mathbb F_9^\times$. List the monic irreducible
polynomials of degrees $1, 2, 3$ over $\mathbb F_2$, and verify
$X^8 - X = X(X+1)(X^3+X+1)(X^3+X^2+1)$ over $\mathbb F_2$.
\end{exercise}

\begin{exercise}[$\star\star$]\label{exo:b3:galois:4}
(a) Find all primitive roots modulo $7$ and modulo $11$ (i.e.\
generators of $\mathbb F_7^\times$, $\mathbb F_{11}^\times$).
(b) Show that for $p$ odd, $x \in \mathbb F_p^\times$ is a
square iff $x^{(p-1)/2} = 1$ (Euler's criterion), and recover
the criterion for $-1$ of \cref{pb:b3:rings:1}.
\end{exercise}

\begin{exercise}[$\star\star$]\label{exo:b3:galois:5}
Show that $\mathbb F_{p^m} \cap \mathbb F_{p^n} = \mathbb
F_{p^{\gcd(m,n)}}$ and $\mathbb F_{p^m}\mathbb F_{p^n} = \mathbb
F_{p^{\operatorname{lcm}(m,n)}}$ inside a fixed algebraic
closure $\bar{\mathbb F}_p$, and describe
$\operatorname{Gal}(\mathbb F_{p^n}/\mathbb F_{p^m})$ for $m
\mid n$.
\end{exercise}

\begin{exercise}[$\star\star$]\label{exo:b3:galois:6}
Let $I_d(q)$ be the number of monic irreducible polynomials of
degree $d$ over $\mathbb F_q$. Prove
\[
X^{q^n} - X \;=\; \prod_{d \mid n}\ \prod_{P \text{ irred.
monic, } \deg P = d} P ,
\qquad\text{hence}\qquad
q^n = \sum_{d \mid n} d\, I_d(q).
\]
Deduce $I_1, I_2, I_3, I_4$ explicitly, and $I_d(q) \geq 1$ for
every $d$ (so extensions $\mathbb F_{q^d}/\mathbb F_q$ exist as
quotients $\mathbb F_q[X]/(P)$ for all $d$).
\end{exercise}

\begin{exercise}[$\star\star$]\label{exo:b3:galois:7}
Determine $\operatorname{Gal}(\Q(\sqrt2,\sqrt3)/\Q)$ and the
complete lattice of intermediate fields. Same question for the
splitting field of $(X^2-2)(X^2-3)(X^2-6)$ --- what do you
notice?
\end{exercise}

\begin{exercise}[$\star\star$]\label{exo:b3:galois:8}
(The cubic, solved by its group) Let $P = X^3 + pX + q \in
\Q[X]$ be irreducible with roots $x_1, x_2, x_3$ and splitting
field $L$. Let $\delta = (x_1 - x_2)(x_1 - x_3)(x_2 - x_3)$ and
$\Delta = \delta^2 = -4p^3 - 27q^2$ (admit this classical
identity or verify it by expanding symmetric functions).
(a) Show $\operatorname{Gal}(L/\Q) \cong A_3$ or $S_3$, according
to whether $\Delta$ is or is not a square in $\Q$.
(b) With $j = \zeta_3$, define the Lagrange resolvents $u = x_1
+ jx_2 + j^2x_3$ and $v = x_1 + j^2x_2 + jx_3$. Show $u^3 + v^3
= -27q$ and $uv = -3p$, and solve for $u^3, v^3$: Cardano's
formulas drop out. Where did solvability of $S_3$ get used?
\end{exercise}

\begin{exercise}[$\star\star$]\label{exo:b3:galois:9}
In $\Q(\zeta_5)$: show that the unique quadratic subfield is
$\Q(\sqrt5)$, via the Gauss sums $\eta_0 = \zeta_5 + \zeta_5^4$,
$\eta_1 = \zeta_5^2 + \zeta_5^3$: compute $\eta_0 + \eta_1$ and
$\eta_0\eta_1$, and deduce $\cos\frac{2\pi}5 = \frac{\sqrt5 -
1}4$. Conclude that the regular pentagon is constructible.
\end{exercise}

\begin{exercise}[$\star\star\star$]\label{exo:b3:galois:10}
Let $K = \mathbb F_p(S, T)$ (rational functions in two
indeterminates) and $L = K(S^{1/p}, T^{1/p})$.
(a) Show $[L:K] = p^2$ and that $\alpha^p \in K$ for every
$\alpha \in L$.
(b) Deduce that $L/K$ is \emph{not} simple: no primitive element
exists --- inseparability is fatal to
\cref{thm:b3:galois:primitive}.
\end{exercise}

\begin{exercise}[$\star\star\star$]\label{exo:b3:galois:11}
Let $P = X^5 - 4X + 2$ and $G$ its Galois group over $\Q$,
acting on the $5$ roots.
(a) Show $P$ is irreducible, and deduce $5 \mid \abs G$; conclude
that $G$ contains a $5$-cycle (Cauchy,
\cref{thm:b3:groups:cauchy}).
(b) Show, by studying the variations of $x \mapsto x^5 - 4x +
2$, that $P$ has exactly $3$ real roots; deduce that complex
conjugation restricts to a transposition in $G$.
(c) Show that a subgroup of $S_5$ containing a transposition and
a $5$-cycle is $S_5$ \emph{(conjugate the transposition by
powers of the cycle)}. Conclude $G \cong S_5$ and, with
\cref{thm:b3:galois:solvable}, that $P$ is not solvable by
radicals.
\end{exercise}

\section{Problem: Gauss and the regular 17-gon}

\begin{problem}[{Weekend problem --- constructibility of the
17-gon}]\label{pb:b3:galois:1}
On March 30, 1796, the nineteen-year-old Gauss showed that the
regular $17$-gon is constructible --- the first progress on the
question since antiquity. We reconstruct his computation with the
tools of this chapter. Set $\zeta = \eu^{2\iu\pi/17}$, $L =
\Q(\zeta)$, $G = \operatorname{Gal}(L/\Q)$.

\medskip
\noindent\textbf{Part I --- The group and its filtration.}
\begin{enumerate}
  \item Justify: $[L:\Q] = 16$, $G \cong (\Z/17\Z)^\times$,
        cyclic of order $16$. Verify that $3$ is a generator of
        $(\Z/17\Z)^\times$ \emph{(compute the powers of $3$
        modulo $17$: $3, 9, 10, 13, 5, 15, 11, 16, \dots$)}.
  \item Let $\sigma \in G$ with $\sigma(\zeta) = \zeta^3$, and
        $H_k = \langle \sigma^{2^k}\rangle$ for $k = 0, \dots,
        4$. Show that $G = H_0 \supset H_1 \supset H_2 \supset
        H_3 \supset H_4 = \{e\}$ with each index $[H_k :
        H_{k+1}] = 2$, and that the fixed fields $\Q = L_0
        \subset L_1 \subset L_2 \subset L_3 \subset L_4 = L$
        form a tower of quadratic extensions.
  \item Conclude \emph{a priori}, using
        \cref{thm:b3:galois:wantzel}, that $\zeta$ --- hence the
        $17$-gon --- is constructible. The rest of the problem
        makes the tower explicit.
\end{enumerate}

\medskip
\noindent\textbf{Part II --- The periods of length 8.}
Define the \emph{Gauss periods}
\[
\eta_0 = \sum_{k \text{ even}} \zeta^{3^k \bmod 17}
= \zeta^{1} + \zeta^{9} + \zeta^{13} + \zeta^{15} + \zeta^{16} +
\zeta^{8} + \zeta^{4} + \zeta^{2},
\qquad
\eta_1 = \sum_{k \text{ odd}} \zeta^{3^k \bmod 17}.
\]
\begin{enumerate}[resume]
  \item Show that $\eta_0, \eta_1$ are fixed by $H_1$ and swapped
        by $\sigma$; deduce $\eta_0, \eta_1 \in L_1$ and that
        they are the two roots of a quadratic over $\Q$.
  \item Compute $\eta_0 + \eta_1 = -1$. Show $\eta_0\eta_1 = -4$
        \emph{(each product $\zeta^a\zeta^b$ is some $\zeta^c$,
        $c \neq 0$; count how many times each $c$ occurs, or
        argue that the product is a rational integer fixed by
        $G$, equal to the sum over all $64$ products, and use
        that each nonzero residue appears equally often)}.
  \item Deduce $\eta_0 = \frac{-1 + \sqrt{17}}2$, $\eta_1 =
        \frac{-1-\sqrt{17}}2$ (identify which is which
        numerically: $\eta_0 \approx 1.56$), and $L_1 =
        \Q(\sqrt{17})$.
\end{enumerate}

\medskip
\noindent\textbf{Part III --- Periods of length 4 and 2.}
Define
\[
\beta_0 = \zeta + \zeta^{13} + \zeta^{16} + \zeta^{4},\quad
\beta_1 = \zeta^3 + \zeta^5 + \zeta^{14} + \zeta^{12},\quad
\beta_2 = \zeta^9 + \zeta^{15} + \zeta^{8} + \zeta^{2},\quad
\beta_3 = \zeta^{10} + \zeta^{11} + \zeta^{7} + \zeta^{6}.
\]
\begin{enumerate}[resume]
  \item Show $\beta_0 + \beta_2 = \eta_0$, $\beta_1 + \beta_3 =
        \eta_1$, and that $\beta_0, \beta_2$ are fixed by $H_2$,
        swapped by $\sigma^2$.
  \item Compute $\beta_0\beta_2 = -1$ and $\beta_1\beta_3 = -1$
        \emph{(expand: the sixteen exponents obtained cover
        $1, \dots, 16$ exactly once)}.
  \item Deduce $\beta_0 = \frac{\eta_0 + \sqrt{\eta_0^2 + 4}}2$
        (check the sign numerically: $\beta_0 \approx 2.05$) and
        the analogous formula for $\beta_1$; hence $L_2 =
        \Q(\beta_0)$, quadratic over $L_1$.
  \item Let $\gamma_0 = \zeta + \zeta^{16} =
        2\cos\frac{2\pi}{17}$ and $\gamma_1 = \zeta^{13} +
        \zeta^4$. Show $\gamma_0 + \gamma_1 = \beta_0$ and
        $\gamma_0\gamma_1 = \beta_1$, so that $\gamma_0 =
        \frac{\beta_0 + \sqrt{\beta_0^2 - 4\beta_1}}2$.
  \item Assemble the chain of formulas expressing
        $\cos\frac{2\pi}{17}$ by nested square roots, and give a
        decimal check ($\cos\frac{2\pi}{17} \approx 0.93247$).
\end{enumerate}

\medskip
\noindent\textbf{Part IV --- Epilogue.}
\begin{enumerate}[resume]
  \item Where exactly did the argument use that $17$ is a
        \emph{Fermat prime} ($17 = 2^{2^2} + 1$)? Show that for
        a prime $p$, the regular $p$-gon is constructible iff $p
        = 2^{2^t} + 1$ for some $t$ \emph{(if $p - 1 = 2^m$,
        show $m$ must itself be a power of $2$)}.
  \item Deduce the complete list of constructible regular
        $n$-gons for $n \leq 20$, using the Gauss--Wantzel
        criterion of \cref{cor:b3:galois:impossible}.
\end{enumerate}
\end{problem}
